Distribution network planning under load and outage uncertainty requires portfolios that remain useful as binding stresses change. SHIELD-MOEA extends non-dominated sorting with selective worst-scenario screening and disjoint scenario draws for search and final evaluation. On a SimBench-derived benchmark with 72 candidate actions, five objectives, and eight archived experiment labels, the method reaches pooled mean hypervolume 0.2740 over 30 seeded runs per stochastic method. The eight labels represent five active proxy configurations because the configured DER-output multiplier does not enter any p4 optimization or scoring objective; the \texttt{der\_uncertainty} block is therefore an inactive-factor replicate, not DER-uncertainty evidence. The pooled result is 5.09\% above NSGA-II with post-hoc final-population repair and 5.56\% above plain NSGA-II. All 32 comparisons with four stochastic baselines are Holm-significant wins; sixteen gaps to two deterministic rules are descriptive and favor SHIELD-MOEA. The margin persists for sampled worst-envelope hypervolume (+5.36\%) and under held-out load/outage ranges (+6.04\%). Feasibility repair carries the main measured gain. Static code-path accounting gives 65\% fewer search-phase plan--scenario objective rows with screening, although no front-quality difference is detected against no screening; hybrid variation and periodic re-screening are also unresolved against simpler controls. A 1296-case composition-level AC check across six networks yields aggregate feasibility 0.685, compared with 0.389 without planning and 0.694 for post-hoc-repaired NSGA-II. The evidence supports a disjoint-evaluation workflow under the implemented load/outage proxies while identifying repair--rather than every named operator--as the resolved performance mechanism.
